\section{Introduction}
\label{sec:intro}

Schistosomiasis remains a major parasitic burden in endemic regions, and
\emph{Schistosoma japonicum} infection drives a distinctive form of hepatic
injury. Eggs trapped in the periportal tracts provoke granulomatous
inflammation and progressive septal fibrosis, producing the reticular,
network-like and tortoise-shell echo patterns that characterize the disease on
B-mode ultrasound. Because ultrasonography is inexpensive, portable and free of
ionizing radiation, it is the practical modality for large-scale screening and
longitudinal follow-up in endemic areas, where transient elastography and
magnetic resonance imaging are rarely
available~\cite{schistoelasto2018,schistopocus2024}.

Ultrasound-based staging, however, depends heavily on operator experience.
Interobserver reproducibility has also been evaluated in shear-wave
elastography for hepatitis C-associated liver fibrosis, with excellent
agreement reported for the point shear-wave method studied by
Kaposi et al.~\cite{interobserver2019}, and a meta-analysis of
two-dimensional shear-wave elastography reported heterogeneity in diagnostic
performance across studies, with differences in sampling protocols and
equipment discussed as possible contributors~\cite{swe2018staging}.
These studies concern elastography rather than automated grading of B-mode
images. A separate body of work applies deep learning to ultrasound-based
fibrosis assessment~\cite{usfibrosisgan2022,microflow2023,hfus2025,liverspleen2025}.

Two properties distinguish the schistosomal grading task from generic fibrosis
classification, and both shape the present work.

First, the clinically meaningful target is \emph{fine-grained}. The
conventional four-grade scale is interpretable but coarse relative to the
continuous progression of septal fibrosis, and a quantized 36-level scale in
steps of $0.1$ better reflects gradual change while remaining mappable to the
familiar clinical grades. Our previous work, SFibAI, established that such
fine-grained continuous scoring is feasible at the image level by predicting a
posterior over the 36 levels and reporting its
expectation~\cite{xu2026sfibai}. Precisely because the target is fine-grained,
the discriminative evidence is subtle: the model must resolve differences in
septal texture and parenchymal architecture that a four-way decision boundary
would never need to distinguish.

Second, acquisition is \emph{anatomically structured}. Images are collected
under a standardized six-position protocol, so each image carries a known
anatomical position, and fibrotic change is annotated with coarse bounding
boxes. Both signals are available at training time, and both are plausibly
informative: the appearance of normal parenchyma differs systematically across
scan positions, so knowing where an image was taken changes what should be
regarded as abnormal.

This raises the question we study. Anatomical position and lesion extent are
readily available as \emph{training} supervision, but they cannot be assumed
present at deployment: position metadata is often missing or unreliable in
routine practice, and box annotation is a research-time artifact. A useful
design must therefore learn from these signals without depending on them at
inference. Whether such anatomical supervision actually helps fine-grained
grading, and whether the two signals help in the same way, is not obvious a
priori.

We introduce \textbf{SynAP-Fib}, which augments the SFibAI grading backbone
with two anatomical prior branches: a six-class position head and a weak-box
lesion attention head. Each branch feeds the grading pathway through a gated
residual injection with a bidirectionally detached gradient boundary, so that
the branches supply anatomical context without becoming additional gradient
routes into the shared representation. At inference the grading pathway
consumes only the model's own \emph{predicted} position posterior and lesion
attention; no position metadata or box annotation is required.

Evaluating the two priors separately and jointly yields the central empirical
observation of this paper. Adding either prior alone degraded the composite
grading risk relative to the baseline, yet supplying both together did not
inherit that degradation and produced the best observed composite risk among
the four configurations. The behavior of the joint configuration is therefore
not what the two single-prior configurations would predict in isolation.

Our contributions are as follows.

\begin{itemize}
\item We present SynAP-Fib, a fine-grained ultrasound liver fibrosis grading
      model in which anatomical position and weak lesion evidence act as
      training-time supervision and as inference-time context through gated,
      gradient-isolated residual injection, requiring no anatomical metadata at
      deployment.
\item We report a $2\times2$ comparison of the two anatomical priors showing a
      non-additive empirical pattern: each prior alone degrades composite
      grading risk, whereas their combination does not, and attains the best
      observed composite risk of the four configurations.
\item We show that the joint configuration supplies two additional clinically
      interpretable outputs -- six-class anatomical position and weak lesion
      localization -- for a $5.9\%$ parameter and $2.8\%$ arithmetic overhead,
      while retaining baseline-level grading accuracy.
\item We evaluate under a prespecified, clinically weighted composite risk that
      combines image-level, patient-level and center-balanced views, with
      checkpoint selection confined to the validation split and the test split
      reserved for reporting.
\end{itemize}

The remainder of the paper is organized as follows. Section~\ref{sec:related}
positions SynAP-Fib against prior work. Section~\ref{sec:methods} describes the
cohort, architecture and evaluation protocol. Section~\ref{sec:results} reports
the results, and Section~\ref{sec:discussion} discusses their interpretation and
limitations.
